\SourceLeaf{352}
\begin{center}\large\bfseries [PROJECTION OF THE CUBE.]\end{center}

\Needspace{8\baselineskip}\begin{center}[1.]\end{center}

If the complex values of the orthographic projections of three straight lines of equal length and perpendicular to one another are
\[
a,\ b,\ c,
\]
then
\[
aa+bb+cc=0,
\]
e.\,g.
\[
a=1+8i,\qquad b=8+i,\qquad c=4-4i.
\]
In general, one can set
\[
a=pp+qq,\qquad b=2ipq,\qquad c=ipp-iqq.
\]

\Needspace{9\baselineskip}\begin{center}[2.]\end{center}

The most elegant form of the solution of the general equation
\[
pp+qq+rr=0
\]
is, with $a$ and $b$ denoting arbitrary complex numbers,
\[
p=(a-b)(b-aii)
\]
\[
q=(b-bi)(aii-ai)
\]
\[
r=(bi-a)(ai-b).
\]

\Needspace{9\baselineskip}\begin{center}\bfseries REMARKS.\end{center}

Part of these two notes, found in a notebook and dating from 1831, has already been printed in vol.\,II, p.\,309.

\indent From December 1839 to Easter 1840, \textsc{Gauss} gave a course of lectures on the \emph{Theorie der imaginären Grössen}, which E.\,\textsc{Heine} attended and wrote up, as mentioned on p.\,330. The following account of the projection of the cube comes from Heine's write-up.
\par\noindent\hfill\textsc{Stäckel}.\par
